% ch02 —— §2 Preliminaries（原文 wcf-paper.tex 559–1164 行）
\section{预备知识}
\label{sec:preliminaries}

\subsection{$d=4$、$\CN=2$ 规范理论}

我们考虑一个 $d=4$、$\CN=2$ 超对称的规范理论：规范群 $G$ 的秩为 $r$，并带有一个特征（复）质量标度 $\Lambda$。\SW 理论（由 \cite{Seiberg:1994rs,Seiberg:1994aj} 开创，更一般的综述见如 \cite{Lerche:1997sm,Freed:1997dp,ias-volumes}）对这类规范理论在其库仑分支上、能量 $\mu \ll \Lambda$ 时的行为给出了相当完整的描述，具体如下。

库仑分支是一个复维数为 $r$ 的复流形 $\vm$（4d 模空间），由矢量多重态标量的真空期望值参数化。我们把 $\vm$ 上的一般坐标系记作 $(u^1, \dots, u^r)$。在每一点 $u \in \vm$，规范群破缺为极大环面 $U(1)^r$。存在一个电磁荷构成的格 $\clat_u \simeq \IZ^{2r}$，其上装备着取整数值的辛配对 $\inprod{,}$。这个格是 $\vm$ 上一个局部系统 $\clat$ 的纤维。也就是说，在 $u \in \vm$ 之上有一个以 $\clat_u$ 为纤维的格的丛，它绕 $\vm$ 中某些奇异轨迹具有非平凡的单值性——那些轨迹的复余维为 $1$，其上有些 BPS 粒子变成无质量。我们有时不严谨地写「$\gamma \in \clat$」，意思是说 $\gamma$ 是 $\clat$ 的一个局部截面。

存在一个「周期」向量 $Z(u) \in \clat_u^* \otimes_\IZ \IC$，它随 $u$ 全纯地变化。对任意 $\gamma \in \clat$，我们如下定义中心荷 $Z_\gamma(u)$：
\begin{equation}
Z_\gamma(u) = Z(u) \cdot \gamma.
\end{equation}
无论是无质量部分还是有质量部分，$Z(u)$ 的描述都起根本的作用。

我们从无质量部分开始。在 $\vm$ 上局部地，可以把 $\clat$ 分裂为 $\clat = \clat^m \oplus \clat^e$，其中两个子格分别由「磁」荷与「电」荷构成，且都是拉格朗日子格。利用 $\clat$ 上的配对，$\clat^m$ 与 $\clat^e$ 互为对偶。这样的分裂称为一个电磁对偶标架。具体地，我们可以为 $\clat^m$ 取一组基 $\{\alpha_1, \dots, \alpha_r\}$，为 $\clat^e$ 取一组基 $\{\beta^1, \dots, \beta^r\}$，使得
\begin{equation}
\inprod{\alpha_I, \alpha_J} = 0, \quad \inprod{\beta^I, \beta^J} = 0, \quad \inprod{\alpha_I, \beta^J} = \delta_I^J
\end{equation}
其中 $I, J = 1, \dots, r$。选定这样的标架后，我们在 $\vm$ 上得到一组「特殊坐标」$a^I$；它们不是别的，正是电中心荷，即
\begin{equation} \label{eq:special-coords}
Z_{\beta^I} = a^I.
\end{equation}
磁中心荷于是是诸 $a^I$ 的全纯函数。它们由单个函数 $\CF(a^I)$（依赖于所选标架）——即 $\CN=2$ 预势——决定：
\begin{equation} \label{eq:prepot}
Z_{\alpha_I} = \frac{\partial \CF}{\partial a^I}.
\end{equation}
这特别蕴含 $Z$ 并非任意：由混合偏导数的对称性可得
\begin{equation} \label{eq:lag-constraint}
\inprod{dZ, dZ} = 0.
\end{equation}
在 \eqref{eq:lag-constraint} 的左端，我们既用了反对称配对 $\inprod{,}$，又用了 $\vm$ 上 $1$-形式的反对称楔积；两者复合而成的配对是对称的，所以这个条件并非空话。事实上，\eqref{eq:lag-constraint} 说的是：在 $u$ 附近，$\CB$ 可以局部地等同于 $\Gamma^*_u \otimes_\IZ \IC$ 的一个复拉格朗日子流形。

预势完全决定了二导数阶的有效拉氏量，后者用电源矢量多重态写出。为写出这个拉氏量，我们引入由
\begin{equation}
\tau_{IJ}(a) = \frac{\partial^2 \CF}{\partial a^I \partial a^J},
\end{equation}
定义的对称矩阵 $\tau$，然后采用一种略去规范指标的记号，例如用 $\tau \abs{da}^2$ 表示 $\tau_{IJ} da^{I} \wedge \star d\bar{a}^J$。于是拉氏量的玻色部分为
\begin{equation} \label{eq:sw-lagrangian}
\CL^{(4)} = \frac{\im \tau}{4\pi} \left(- \abs{da}^2 - F^2 \right) + \frac{\re \tau}{4\pi} F \wedge F.
\end{equation}

中心荷 $Z_\gamma$ 对有质量谱同样具有根本的重要性。事实上，任何电荷为 $\gamma$ 的单粒子态，其质量都满足
\begin{equation} \label{eq:bps-bound}
M \ge \abs{Z_\gamma}
\end{equation}
且等号成立当且仅当该态是 BPS 态。BPS 态属于超庞加莱对称的有质量短多重态；在小群 $SU(2)$ 下，这样一个多重态中的静止态变换为
\begin{equation}
[j] \otimes \left(\left[1/2\right] + 2[0]\right).
\end{equation}
取 $j=0$ 给出有质量超多重态，而 $j = \half$ 则是有质量矢量多重态。

有一个标准的指标可以「数」短多重态，即第二螺旋度超迹 $\degen(\gamma;u)$。对 $u \in \vm$ 处理论谱中每一个电荷为 $\gamma$ 的有质量超多重态，这个超迹贡献 $+1$；类似地，对每一个有质量矢量多重态贡献 $-2$。在任何使单粒子态不与多粒子态连续谱混合的形变下，$\degen(\gamma;u)$ 不变。由 \eqref{eq:bps-bound} 可知，这种混合受到很强的限制：一个 BPS 粒子只能衰变为其他 BPS 粒子，而且只有当它们的中心荷全都具有同一相位时才行。因此，在 $\vm$ 中远离\emph{边缘稳定墙}处，$\degen(\gamma;u)$ 关于 $u$ 局部常值。这些边缘稳定墙的实余维为 $1$；对一对线性无关的电荷 $\gamma, \gamma'$，它们定义为使 $Z_\gamma$ 与 $Z_{\gamma'}$ 非零且具有同一相位的诸 $u\in \vm$ 所构成的轨迹。

理解 $\degen(\gamma;u)$ 在 $u$ 穿过边缘稳定墙时的跳变行为，是本文的主要动机之一。下面我们就转向它。

\subsection{Kontsevich--Soibelman 穿墙公式}
\label{sec:ReviewKS}

本节回顾 Kontsevich--Soibelman 穿墙公式。如 \cite{ks1} 最初所提出的，该公式决定了「广义 Donaldson--Thomas 不变量」$\hat\degen(\gamma;u)$ 的跳变行为。下文将会看到，若把 Donaldson--Thomas 不变量与螺旋度超迹等同起来，即 $\hat\degen(\gamma;u) = \degen(\gamma;u)$，那么穿墙公式在若干非平凡的例子中给出物理上所预期的答案：特别地，它重现了 \cite{Denef:2007vg} 的「原始穿墙公式」，也重现了规范群为 $SU(2)$ 的 \SW 理论之 BPS 谱的穿墙行为。

有一个技术要点：为使 KS 公式有意义，诸 $\degen(\gamma;u)$ 不能完全任意。在 $\clat$ 上引入一个正定范数后，必须要求存在某个 $K > 0$，使得
\begin{equation}
\frac{\abs{Z_\gamma}}{\norm{\gamma}} > K
\end{equation}
对所有满足 $\hat\degen(\gamma;u) \neq 0$ 的 $\gamma$ 成立。本文通篇假定这条性质——称为「支撑性质」（Support Property）——成立。

\subsubsection*{Kontsevich--Soibelman 代数}

穿墙公式是用一个李代数给出的；该李代数由生成元 $e_\gamma$（$\gamma \in \clat$）以及一条基本对易关系定义：
\begin{equation} \label{eq:ks-commutator}
[e_{\gamma_1}, e_{\gamma_2}] = (-1)^{\langle
\gamma_1,\gamma_2\rangle} \langle \gamma_1,\gamma_2\rangle
e_{\gamma_1+\gamma_2}.
\end{equation}
本文将把这个抽象的李代数实现为一个复化环面上的复辛同胚代数，这一点很重要。除去本节末将出现的一个细微之处，该环面就是局部系统 $\tilde{T} := \clat^* \otimes_{\IZ} \IC^\times$ 的纤维 $\tilde{T}_u$。

任意 $\gamma \in \clat$ 给出 $\tilde{T}_u$ 上一个相应的函数 $X_\gamma$，满足 $X_\gamma X_{\gamma'} = X_{\gamma + \gamma'}$。取定 $\clat$ 的一组基 $\{\gamma^1, \dots, \gamma^{2r}\}$ 后，我们可以取 $X^i := X_{\gamma^i}$ 作为 $\tilde{T}_u$ 的坐标。$\clat^*$ 上的辛配对给出 $\tilde{T}_u$ 上一个全纯辛形式 $\hsymp^{\tilde{T}}$：若 $\epsilon^{ij} = \inprod{\gamma^i, \gamma^j}$，而 $\epsilon_{ij}$ 是它的逆，则
\begin{equation}
\hsymp^{\tilde{T}} = \frac{1}{2}\epsilon_{ij} \frac{dX^i}{X^i} \wedge \frac{dX^j}{X^j}.
\end{equation}

我们希望把 $e_\gamma$ 等同于 $\tilde{T}_u$ 上由哈密顿量 $X_\gamma$ 生成的无穷小辛同胚。这差不多给出代数 \eqref{eq:ks-commutator}，但漏掉了额外的符号 $(-1)^{\langle \gamma_1,\gamma_2\rangle}$。这个符号在下文与物理上已知的诸穿墙公式作比较时很重要；在那种语境下，它与下面的事实有关：由两个电荷为 $\gamma_1,\gamma_2$ 的粒子组成的束缚态，其费米子数会平移 $\langle \gamma_1,\gamma_2\rangle$。

在 $\vm$ 的一个局部片上，我们可以通过引入 $\IZ_2$ 值二次型 $(-1)^{\langle \gamma_1,\gamma_2\rangle}$ 的一个「二次加细」（quadratic refinement）来吸收这个符号：所谓二次加细，是指一个满足
\begin{equation}
\sigma(\gamma_1)\sigma(\gamma_2) = (-1)^{\langle \gamma_1,\gamma_2\rangle} \sigma(\gamma_1 + \gamma_2)
\end{equation}
的映射 $\sigma: \clat \to \IZ_2$。得到这样的 $\sigma$ 的一个办法是：取一个局部电磁对偶标架 $\clat \cong \clat^e \oplus \clat^m$，写 $\gamma = \gamma^e + \gamma^m$，并置 $\sigma(\gamma)=(-1)^{\inprod{\gamma^e,\gamma^m}}$。注意
\begin{equation}
\sigma(\gamma_1)\sigma(\gamma_2) = \sigma(\gamma_1 + \gamma_2)
(-1)^{\inprod{\gamma_1^e,\gamma_2^m} + \inprod{\gamma_2^e,\gamma_1^m}} =
\sigma(\gamma_1 + \gamma_2) (-1)^{\langle \gamma_1,\gamma_2\rangle}
\end{equation}
正如所需。无论如何，选定了某个 $\sigma(\gamma)$ 之后，我们便可以把 $e_\gamma$ 等同于由哈密顿量 $\sigma(\gamma) X_\gamma$ 生成的辛同胚。

任意两个加细 $\sigma$、$\sigma'$ 都满足 $\sigma(\gamma) \sigma'(\gamma) = (-1)^{c(\sigma, \sigma') \cdot \gamma}$，其中 $c(\sigma, \sigma') \in \clat^* / 2 \clat^*$ 是某个固定的元素。与这两个加细相联系的哈密顿量 $\sigma(\gamma) X_\gamma$ 与 $\sigma'(\gamma) X_\gamma$，则由 $\tilde{T}_u$ 的一个自同构相联系，该自同构把 $X_\gamma$ 送到 $(-1)^{c(\sigma, \sigma') \cdot \gamma} X_\gamma$。

\subsubsection*{穿墙公式}

现在我们准备表述穿墙公式。它的基本构件是群元素
\begin{equation}
\CK_\gamma := \exp \sum_{n=1}^{\infty} \frac{1}{n^2} e_{n \gamma}.
\end{equation}
在我们的等同之下，$\CK_\gamma$ 成为一个作用在 $\tilde{T}_u$ 上的辛同胚，由
\begin{equation} \label{eq:ks-symplect}
\CK_\gamma : X_{\gamma'} \to X_{\gamma'} (1- \sigma(\gamma)
X_\gamma)^{\langle \gamma',\gamma \rangle}
\end{equation}
给出。

对每个电荷为 $\gamma$ 的 BPS 粒子，联系一条复 $\zeta$ 平面中的射线，它由中心荷决定：
\begin{equation}
\ell_\gamma := \{\zeta: Z_\gamma(u) / \zeta \in \IR_- \}.
\end{equation}
当 $u \in \vm$ 变动时，这些射线在 $\zeta$ 平面中转动。射线的循环次序只在 $u$ 到达边缘稳定墙时才改变。在这样一面墙上，一组 BPS 射线 $\ell_\gamma$ 汇聚到一起，它们对应一组使诸 $Z_\gamma$ 变得共线的电荷。在边缘稳定墙上的一般点处，这组电荷可以参数化为\footnote{为确立这些 $\gamma_1$、$\gamma_2$ 的存在性，我们需要用到支撑性质：否则很容易设想诸共线的 $Z_\gamma$ 在原点附近堆积的情形。} $\{n \gamma_1 + m \gamma_2: m, n > 0\}$，其中 $\gamma_1, \gamma_2$ 是某些原始向量，且 $Z_{\gamma_1} / Z_{\gamma_2} \in \IR_+$。

现在，对每个电荷为 $\gamma$ 的 BPS 粒子联系群元素 $\CK_\gamma$，并对在墙上变得共线的诸态取乘积：
\begin{equation} \label{eq:ks-product}
A := \prod^\cwarrow_{\substack{\gamma = n \gamma_1 + m \gamma_2 \\ m>0,\,n>0}}
\CK_\gamma^{\degen(\gamma;u)},
\end{equation}
其中各因子的次序对应于诸射线 $\ell_\gamma$ 的顺时针次序。我们可以对墙两侧的 $u$ 分别考虑这个乘积。当 $u$ 穿过墙时，因子的次序反转，而诸 $\degen(\gamma;u)$ 跳变。穿墙公式的陈述是：整个乘积 $A$ 不变。

这个条件强得足以由诸 $\degen(\gamma;u_-)$ 定出诸 $\degen(\gamma;u_+)$，其中 $u_\pm$ 是在墙的两侧作无穷小位移所得的点。为理解实践中如何做到这一点，我们首先得处理一个重要而细微之处：由于 BPS 态的谱通常是无限的，乘积 \eqref{eq:ks-product} 一般涉及无穷多个因子。按照 \cite{ks1}，我们可以这样理解它。该乘积只涉及诸生成元 $e_{n \gamma_1 + m \gamma_2}$，其中 $m,n>0$。它们生成的李代数可以通过固定某个整数 $L$、并令 $e_{n \gamma_1 + m \gamma_2} = 0$（只要 $n+m > L$）来一致地截断。（也就是说，该李代数被一族满足 $n+m>L$ 的李子代数过滤，而我们可以依次模去 $L$ 越来越大的那些子代数。）经过这样的截断，\eqref{eq:ks-product} 只涉及有限多个非平凡项；无穷乘积则理解为这些截断乘积当 $L \to \infty$ 时的极限。

本着同样的精神，考虑辛同胚 $A$ 的幂级数展开
\begin{equation}
A : X_{\gamma'} \to  ( 1 + \sum_{m>0,n>0} c^{m,n}_{\gamma'} X_{n \gamma_1 + m \gamma_2}) X_{\gamma'},
\end{equation}
并将它截断到 $n+m \le L$。我们可以在边缘稳定墙的一侧算出这一展开，然后在墙的另一侧递归地定出诸 $\degen(\gamma;u)$。具体地说，先取 $L=1$；此时 $\degen(\gamma_1;u)$ 与 $\degen(\gamma_2;u)$ 由「正确地重现 $c^{1,0}_{\gamma'}$ 与 $c^{0,1}_{\gamma'}$」这一要求定出。接着取 $L=2$，考虑 $A \CK_{\gamma_1}^{-\degen(\gamma_1;u)} \CK_{\gamma_2}^{-\degen(\gamma_2;u)}$ 的展开，以提取下一组简并度。这一迭代可以直接地继续下去，定出所有的 $\degen(n \gamma_1 + m \gamma_2 ; u)$。远非显然——但 \cite{ks1} 已作猜想——的是：这样算出的诸 $\degen(\gamma ; u)$ 都是整数！

\subsubsection*{几个例子}

在上述对 Kontsevich--Soibelman 公式的诠释中，我们把他们的广义 Donaldson--Thomas 不变量与物理上定义的 $\degen(\gamma;u)$ 等同了起来。为给这一等同提供动机，下面我们描述几个例子。

如上所述，在边缘稳定墙的一般点上，进入 WCF 的诸辛同胚由一个二维的电荷格生成，即 $\gamma = (p,q) \in \IZ^2$，其典范辛形式为 $\inprod{(p,q), (p',q')} = pq' - q p'$。相应地，我们写 $X_{1,0}=x$，$X_{0,1}=y$。于是辛同胚 $\CK_{p,q}$ 由它们在 $x$ 与 $y$ 上的作用完全确定，具体为
\begin{equation}
\CK_{p,q}: (x,y) \to \biggl(  \bigl(1-(-1)^{pq} x^p y^q \bigr)^q x,
\bigl(1-(-1)^{pq} x^p y^q\bigr)^{-p} y \biggr).
\end{equation}

考虑这样一面边缘稳定墙：一个原始电荷为 $(1,0)$ 的单个 BPS 粒子与一个原始电荷为 $(0,1)$ 的单个粒子，其中心荷在这面墙上汇合。Kontsevich 与 Soibelman 注意到一个漂亮的「五边形恒等式」：
\begin{equation} \label{eq:product-formula-primitive}
\CK_{1,0} \CK_{0,1} = \CK_{0,1} \CK_{1,1} \CK_{1,0}.
\end{equation}
于是 WCF 预言：穿过该墙时只会多产生一个粒子，即由一个荷电粒子与一个荷磁粒子组成的双荷束缚态。确实，来自超引力的「原始穿墙公式」（它在场论中也成立）\cite{Denef:2007vg} 预言：这对粒子会形成一个处于超多重态表示中的单一束缚态。它还预言：一个电荷为 $(1,0)$ 的单个粒子不能与多于一个电荷为 $(0,1)$ 的粒子束缚在一起。要数清由若干个不同类型的粒子组成的更一般的束缚态是相当困难的。它们的不存在，已经是 KS 穿墙公式的一个非平凡预言。

与原始穿墙公式的进一步比较，有助于我们理解 $e_\gamma$ 的对易关系 \eqref{eq:ks-commutator} 中那个符号的作用。考虑乘积 $\CK_{\gamma_1} \CK_{\gamma_2}$，试着把它改写成沿相反方向排列（即让诸 $Z_{\gamma_i}$ 的斜率递增而非递减）的乘积，形如 $\CK_{\gamma_2} \cdots \CK_{\gamma_1}$。设 $\gamma_1,\gamma_2$ 都是原始向量，并考虑由诸 $e_{n\gamma_1 + m \gamma_2}$ 生成的子代数，再模去其中 $n\geq 2, m\geq 2$ 的部分。结果是一个海森堡代数。截断后的海森堡群中的 KS 公式为
\begin{equation}
\CK_{\gamma_1}^{\degen(\gamma_1;u_+)}\CK_{\gamma_1+\gamma_2}^{\degen(\gamma_1+\gamma_2;u_+)}
\CK_{\gamma_2}^{\degen(\gamma_2;u_+)}=
\CK_{\gamma_2}^{\degen(\gamma_2;u_-)}\CK_{\gamma_1+\gamma_2}^{\degen(\gamma_1+\gamma_2;u_-)}
\CK_{\gamma_1}^{\degen(\gamma_1;u_-)}
\end{equation}
其中 $u_\pm$ 是在墙两侧作无穷小位移所得的点。现在，在边缘稳定墙上的一般点处，对 $i=1,2$ 有 $\degen(\gamma_i; u_+) = \degen(\gamma_i; u_-)$。此外，$\CK_{\gamma_1+\gamma_2}$ 在海森堡群中是中心的，因此计算群对易子就重现了原始穿墙公式的如下推论：
\begin{equation}
\Delta \degen = (-1)^{\langle \gamma_1,\gamma_2\rangle -1} \langle
\gamma_1,\gamma_2\rangle \degen(\gamma_1;u) \degen(\gamma_2;u).
\end{equation}
这一论证的一个更精细的版本，可以从 KS 公式提取出 \cite{Denef:2007vg} 的半原始穿墙公式。\footnote{这是与 Wu-yen Chuang 合作的未发表工作中证明的。}

\eqref{eq:product-formula-primitive} 中的例子是特殊的：它的两端都只含有限多项。更典型的情形是遇到无穷乘积。Kontsevich 与 Soibelman 给出的第二个漂亮例子是
\begin{equation}
\CK_{1,0}^2 \CK_{0,1}^2 = \left( \CK_{0,1}^2 \CK_{1,2}^2 \CK_{2,3}^2 \cdots \right)
\CK_{1,1}^{4} \CK_{2,2}^{-2} \left( \cdots \CK_{3,2}^2 \CK_{2,1}^2 \CK_{1,0}^2 \right).
\end{equation}
我们在附录 A 中给出这个恒等式的一个有启发性的证明。% TODO 合稿改 \ref{app:ks-product}
通过一次基变换，我们得到一个物理上很有意思的公式，
\begin{equation} \label{eq:nf2}
\CK_{1,-1}^2 \CK_{0,1}^2 = \left( \CK_{0,1}^2 \CK_{1,1}^2 \CK_{2,1}^2 \cdots \right) \CK_{1,0}^4 \CK_{2,0}^{-2}
\left( \cdots \CK_{3,-1}^2 \CK_{2,-1}^2 \CK_{1,-1}^2 \right),
\end{equation}
它刻画了带两个无质量味的 $SU(2)$ Seiberg--Witten 理论的谱（更确切地说，是超多重态在 $SO(4)=SU(2)_A \times SU(2)_B$ 味对称的矢量表示下变换的情形），如 \cite{Bilal:1996sk} 所述。\footnote{恒等式 \eqref{eq:nf2} 与 Seiberg--Witten 理论的关系，最早由 Frederik Denef 提示；\eqref{eq:nf2} 与 $N_f=2$ 理论的确切关系，则是与 Wu-yen Chuang 合作搞清的。}在右端我们看到完整的弱耦合谱：一个电荷为 $(2,0)$ 的 $W$ 玻色子（它对螺旋度超迹贡献 $-2$）、四个电荷为 $(1,0)$ 的超多重态，以及一组电荷为 $(n,\pm 1)$、重数为 $2$ 的双荷子。（事实上，这些双荷子处于 $SU(2)_A$ 或 $SU(2)_B$ 的双重态中，视 $n$ 的奇偶性而定。）在左端我们看到强耦合谱：一个重数为 $2$ 的单极子（$SU(2)_A$ 的双重态），以及一个重数为 $2$ 的双荷子（$SU(2)_B$ 的双重态）。

变量的一次小小改变 $y \to -y^2$ 把乘积公式 \eqref{eq:nf2} 变成
\begin{equation} \label{eq:nf0}
\CK_{2,-1} \CK_{0,1} = \left( \CK_{0,1} \CK_{2,1} \CK_{4,1} \cdots \right)
\CK_{2,0}^{-2} \left( \cdots \CK_{6,-1} \CK_{4,-1} \CK_{2,-1} \right).
\end{equation}
这个公式刻画了纯 $SU(2)$ Seiberg--Witten 理论的穿墙行为。\footnote{\eqref{eq:nf2} 与 \eqref{eq:nf0} 之所以极为相似，是因为零质量 $N_f=2$ 理论的 SW 曲线是 $N_f=0$ 理论 SW 曲线的二重覆盖。}左端包含强耦合处存在的两个粒子 \cite{Ferrari:1996sv}：一个电荷为 $(0,1)$ 的单极子与一个电荷为 $(2,-1)$ 的双荷子。右端包含弱耦合处的无穷多个双荷子，以及 $W$ 玻色子的贡献 $\CK_{2,0}^{-2}$。

\subsubsection*{加入味信息}

乘积 \eqref{eq:nf2} 描述了 $N_f = 2$ 的 $SU(2)$ Seiberg--Witten 理论的 BPS 谱，但不携带 BPS 粒子的味荷信息。下面我们描述 KS 公式的一个猜想性变体，它把味荷的信息包含进来。（这一公式的物理动机将在第 6 节看到。）% TODO 合稿改 \ref{sec:masses}
引入一个新的格，即味荷格 $\clat^f$，以及一个新参数 $\log \mu \in (\clat^f)^* \otimes_{\IZ} \IC^\times$。然后把诸 $X_\gamma$ 推广为以 $(\gamma,\gamma^f) \in \clat \oplus \clat^f$ 标记的新函数：\footnote{严格说来，完整的电荷局部系统 $\hat\Gamma$ 整体上并不分裂为 $\Gamma \oplus \Gamma^f$；我们实际上有一个扩张 $0 \to \Gamma^f \to \hat\Gamma \to \Gamma \to 0$。不过，这个扩张总可以局部地分裂，而这对我们的目的已经足够。}令 $a$ 取遍 $\clat^f$ 的一组基，
\begin{equation}
X_{\gamma,\gamma^f} := \prod_i (X^i)^{\gamma_i} \prod_a
(\mu^a)^{\gamma^f_a} = X_\gamma \prod_a (\mu^a)^{\gamma^f_a}.
\end{equation}
再定义携带味信息的加细辛同胚：
\begin{equation}
\CK_{\gamma,\gamma^f} : X_{\gamma'} \to X_{\gamma'} (1- \sigma(\gamma)
X_{\gamma,\gamma^f})^{\langle \gamma',\gamma \rangle}.
\end{equation}
中心荷现在依赖于诸质量 $m^a$，即 $Z_{\gamma,\gamma^f}(u) = Z_\gamma(u) + \gamma^f_a m^a$，并确定新的边缘稳定墙。（诸 $\mu^a$ 是诸 $m^a$ 的函数。见下文第 6 节。）% TODO 合稿改 \ref{sec:masses}
我们引入一个与 \eqref{eq:ks-product} 类似的乘积，
\begin{equation} \label{eq:ks-product-masses}
A := \prod^\cwarrow_{\substack{(\gamma,\gamma^f) = n \gamma_1 + m \gamma_2 \\ m>0,\,n>0}}
\CK_{\gamma,\gamma^f}^{\degen(\gamma,\gamma^f;u)},
\end{equation}
推广后的 WCF 断言：$A$ 穿过墙时连续。我们在附录 A 中导出无穷乘积 \eqref{eq:nf2} 的一个含味荷的加细版本；把它与推广后的 WCF 结合起来，就得到 $N_f = 2$ 的 $SU(2)$ 理论正确的穿墙行为。% TODO 合稿改 \ref{app:ks-product}

\subsubsection*{整体性问题}

到目前为止，本节都在 $\vm$ 的一个局部片上工作，并固定了一个二次加细 $\sigma$，以便把 Kontsevich--Soibelman 代数等同于复化环面 $\tilde{T}_u$ 的辛同胚代数——该环面是局部系统 $\tilde{T}$ 的一个纤维。一般而言，不可能在 $\vm$ 上整体地选取这样的加细，因为局部系统 $\clat$ 有单值性。因此，「KS 代数就是作用在 $\tilde{T}$ 上的辛同胚代数」这一说法整体上并不成立。

不过，通过对 $\tilde{T}$ 作适当的扭曲，我们可以定义一个与它密切相关的复化环面丛 $T$，KS 代数确实作用在它上面。$T$ 的定义使得：局部选取一个二次加细就给出一个等同 $T \simeq \tilde{T}$；而给定两个不同的加细 $\sigma, \sigma'$，相应的两个等同相差 $\tilde{T}$ 上的映射 $X_\gamma \to (-1)^{c(\sigma, \sigma') \cdot \gamma} X_\gamma$。纤维方向的辛形式 $\hsymp^{\tilde T}$ 诱导出 $T$ 上相应的纤维方向辛形式 $\hsymp^T$。

具有上述性质的扭曲丛 $T$ 可以这样构造。设 $R$ 记 $\vm$ 上以诸加细 $\sigma$ 为局部截面的那个局部系统。$R$ 是 $\clat^*/2\clat^*$ 的一个挠子（torsor），而 $T$ 是与之相伴的丛：
\begin{equation}
T := \left( \tilde{T} \times R \right) / \left( (X_\gamma, \sigma) \sim
((-1)^{c(\sigma, \sigma') \cdot \gamma} X_\gamma, \sigma') \right).
\end{equation}

\subsection{$\IR^3 \times S^1$ 上的低能有效理论}
\label{sec:reduction}

我们的目标是把 Kontsevich--Soibelman 的 WCF 解释为一个关于 $\IR^3 \times S^1$ 上规范理论的陈述，其中 $S^1$ 的半径为 $R$。我们在一个比其他一切标度都低的能标 $\mu$ 处研究该理论，即 $\mu \ll \Lambda$，且 $\mu \ll 1/R$。在这个能量下，理论看起来有效地是三维的。本节描述它的一些基本性质。

在大半径极限 $R \gg 1/\Lambda$ 下，我们可以用红外拉氏量 \eqref{eq:sw-lagrangian} 确定三维动力学。动力学自由度不过是四维场中不依赖于 $x^4$ 的那些模式。它们当然包括诸标量 $a^I$。此外，从规范扇区我们得到「电」Wilson 线
\begin{equation}
\theta_\elec^I := \oint_{S^1} A_4^I dx^4,
\end{equation}
以及另一组周期标量 $\theta_{\magn,I}$，它们由对偶化 $d=3$ 规范场 $A^I_\alpha dx^\alpha$ 得到。我们常常把它们看作「磁」Wilson 线，
\begin{equation}
\theta_{\magn,I} := \oint_{S^1} (A_{D,4})_I dx^4.
\end{equation}
我们可以用两种方式定义 $\theta_{\magn,I}$：或者在一个把规范场当作自对偶的表述中工作，或者在固定的磁量子数 $P^I$ 下工作、把 $\theta_{\magn,I}$ 作为它们的傅里叶对偶引入。

所有这些周期标量在任一固定的 $u \in \vm$ 处构成一个 $2r$-环面 $\CM_u$ 的坐标。让 $u$ 变动，我们就得到一个环面丛 $\CM$（3d 模空间）。纤维 $\CM_u$ 在 $\vm$ 的奇异轨迹之上退化。$\IR^3$ 上的低能理论是一个以 $\CM$ 为靶空间的 sigma 模型。

更确切地说，$\theta = (\theta_\elec^I, \theta_{\magn,I})$ 是 $2r$-环面的一个局部系统的纤维中的元素，该局部系统记作 $\widetilde\CM := \clat^*\otimes_{\IZ} (\IR/2\pi \IZ)$。对任意 $\gamma \in \clat$，我们得到 $\widetilde\CM$ 上一个角坐标，记作 $\theta_\gamma := \gamma \cdot \theta$。$\widetilde\CM$ 与 $\CM$ 并不完全相同；这里存在一个我们上文略过的整体扭曲，本节末将讨论它。

总之，三维理论是一个映入黎曼流形 $\CM$ 的 sigma 模型；$\CM$ 的实维数为 $4r$，拓扑上是 $\vm$ 上的 $2r$-环面丛。该理论具有 $\CN=4$ 超对称（$8$ 个实超荷），这蕴含 $\CM$ 上的度量是超凯勒的。这个度量是本文研究的主要对象。它此前已在 \cite{Seiberg:1996nz} 中研究过，那里特别把纯 $SU(2)$ 规范理论的 $R \to 0$ 极限识别为 Atiyah--Hitchin 流形。本文更关心相反的极限 $R \to \infty$，因为在这个极限下可以读出 $d=4$ 理论全部 BPS 谱的印记。下一节我们先考虑该极限下的首项行为。

\subsubsection*{整体性问题}

在上面的描述中，我们对 Wilson 线的确切定义稍嫌天真。我们的描述在 $\vm$ 的一个局部片上够用，但正如我们将在第 4 节看到的，它在整体上不可能完全正确。% TODO 合稿改 \ref{sec:periodicnut}
事实上，为使 $\CM$ 上的度量光滑，我们将会看到：绕 $\vm$ 中路径的单值性一般必须伴随 Wilson 线平移 $\pi$。这与我们的天真描述相矛盾，因为环面丛 $\widetilde\CM$ 自带一个零截面。

我们提出，正确的整体图像如下：在任一固定的 $u \in \vm$ 处，Wilson 线位于一个环面 $\CM_u$ 之中，它与 $\widetilde\CM_u$ 同构，但并非典范同构。选定 $\clat_u$ 上二次型 $(-1)^{\inprod{\gamma_1,\gamma_2}}$ 的一个加细 $\sigma$ 后，就得到一个同构 $\CM_u \simeq \widetilde\CM_u$。\footnote{这类二次加细在自对偶规范理论的精确表述中频繁出现 \cite{Witten:1996hc,Freed:2000ta,Hopkins:2002rd,Belov:2006jd}。这里的 $\sigma$ 的来源，很可能可以用这种方式解释。}这样的加细一般只局部存在，所以丛 $\CM$ 与 $\widetilde\CM$ 整体上并不相同。给定两个局部加细 $\sigma$、$\sigma'$，相应的两个局部同构 $\CM \simeq \widetilde\CM$ 相差 $\widetilde\CM$ 上的平移 $\theta \to \theta + \pi c(\sigma,\sigma')$。当然，这一讨论与我们在第 \ref{sec:ReviewKS} 节末所述环面丛 $T$ 与 $\tilde{T}$ 之间的关系是密切平行的。

\subsection{半平坦几何}
\label{sec:semiflat}

在 $R \to \infty$ 极限下，$\CM$ 上度量的首项行为由 $d=3$ 有效作用量支配，后者通过把 \eqref{eq:sw-lagrangian} 简单地截断到不依赖于 $x^4$ 的扇区而得到。这给出
\begin{equation}
\CL^{(3)} = (\im \tau) \left(- \frac{R}{2} \abs{da}^2 - \frac{R}{2}
F^{(3)} \wedge \star F^{(3)} -  \frac{1}{8 \pi^2 R} d\theta_\elec^2
\right) +  (\re \tau) \left(\frac{1}{2\pi} d \theta_\elec \wedge
F^{(3)}\right).
\end{equation}
然后，把 $d=3$ 规范场 $A^I$ 对偶化为标量 $\theta_{\magn,I}$，经过少许整理便得
\begin{equation}
\CL^{(3)}_{dual} = - \frac{R}{2} (\im \tau) \abs{da}^2 - \frac{1}{8 \pi^2 R} (\im \tau)^{-1} \abs{d \theta_\magn - \tau d\theta_\elec}^2.
\end{equation}
这是映入 $\CM$ 的 sigma 模型的拉氏量，其度量局部地由下式给出
\begin{equation} \label{eq:semiflatmetric}
g^\sfm = R (\im \tau) \abs{da}^2 + \frac{1}{4 \pi^2 R} (\im \tau)^{-1} \abs{dz}^2,
\end{equation}
其中我们引入了
\begin{equation}
dz_I = d \theta_{\magn,I} - \tau_{IJ} d \theta_\elec^J.
\end{equation}
（这个记号非常方便，但我们要强调：「$dz_I$」这种形式在整个 $\CM$ 上并不闭合：它只在限制到每个环面纤维 $\CM_u$ 上时才闭合。）

我们称 $g^\sfm$ 为 $\CM$ 上的「半平坦」度量，因为在这个度量下环面纤维是平坦的。表达式 \eqref{eq:semiflatmetric} 反映了这样一个事实：$g^\sfm$ 关于 $\CM$ 上的某个复结构是凯勒的，而在该复结构下 $da^I$ 与 $dz_I$ 构成 $\Omega^{1,0}$ 的一组基。在这个复结构下，$\CM$ 就是 $\vm$ 上由紧\emph{复}环面构成的 Seiberg--Witten 纤维化。（我们将把它与 $\CM$ 上的其他复结构相对照，那些复结构稍后就会遇到；在那些复结构下，诸环面 $\CM_u$ \emph{不是}复子流形。）诸纤维 $\CM_u$ 的体积都是
\begin{equation}
\vol(\CM_u) = \left(\frac{1}{R}\right)^r.
\end{equation}

表达式 \eqref{eq:semiflatmetric} 只在局部有效，因为它用到了一个对偶标架的选择。尽管如此，不同标架下的表达式仍然拼合成一个光滑的度量——除了在 $\vm$ 的奇异轨迹之上，那里 $g^\sfm$ 有奇点。从有效场论的观点看，这样的奇点本不该出现；我们将会看到，它在精确的量子修正度量 $g$ 中由 BPS 瞬子修正所消解。

\begin{yizhu}
$\tilde{T}$ 与 $T$、$\widetilde\CM$ 与 $\CM$ 是两对不同的对象，后文会反复用到，读者务必分清。带波浪号者（$\tilde{T}$、$\widetilde\CM$）是局部系统，只在 $\vm$ 的局部片上才有定义；不带波浪号者（$T$、$\CM$）是整体定义的扭曲丛。二者之间的等同依赖二次加细 $\sigma$ 的局部选取：换一个加细，等同就差一个平移 $X_\gamma \to (-1)^{c(\sigma,\sigma')\cdot\gamma}X_\gamma$（在环面情形为 $\theta \to \theta+\pi c(\sigma,\sigma')$）。本块 §2.2 末与 §2.3 末的两段「整体性问题」讨论的正是这件事。
\end{yizhu}
